\section{Results}
\label{sec:results}

\subsection{FP32 PC$\leftrightarrow$board parity (measured on-board)}

In FP32, the prediction-by-prediction parity is \emph{exact} in the frozen protocol (\num{1.000}) on both
datasets, and near-exact online (\num{0.9754} Pronostia, \num{0.9887} Monitoring) --- the online gaps stem
from the board's \texttt{float32} versus the PC's \texttt{float64} at decision boundaries.
Table~\ref{tab:parity} summarizes; Figure~\ref{fig:parity} illustrates.

\begin{table}[t]
\centering
\caption{FP32 PC$\leftrightarrow$board parity and performance (\emph{measured on-board}, \texttt{exp\_S36}).}
\label{tab:parity}
\begin{tabular}{lccccc}
\toprule
Dataset & Cond. & F1 (PC/board) & Frozen parity & Online parity \\
\midrule
Pronostia  & 5feat & \num{0.9164} & \num{1.000} & \num{0.9754} \\
Pronostia  & all   & \num{0.9180} & \num{1.000} & \num{0.9626} \\
Monitoring & 5feat & \num{0.9194} & \num{1.000} & \num{0.9887} \\
\bottomrule
\end{tabular}
\end{table}

\begin{figure}[t]
\centering
\includegraphics[width=0.9\linewidth]{fig1_parity_fp32_pc_board.png}
\caption{FP32 PC$\leftrightarrow$board parity: frozen (\num{1.000}) vs online. \emph{Measured on-board}
(\texttt{exp\_S36\_parity\_*.json}).}
\label{fig:parity}
\end{figure}

\paragraph{Latency and RAM (Gap~2/Gap~3).} All board latencies stay three orders of magnitude below budget:
inference \SI{50}{\micro\second} (Pronostia) / \SI{48}{\micro\second} (Monitoring); inference + online update
\SI{251}{\micro\second} / \SI{239}{\micro\second}, versus $\approx$\,\SIrange{33}{35}{\micro\second} for the
reference PC inference. The \texttt{.bss} footprint ranges over \SIrange{100}{145}{\kilo\byte}, below the
\SI{256}{\kilo\byte} available. Figure~\ref{fig:latency} places these values against the
\SI{100}{\milli\second} line.

\begin{figure}[t]
\centering
\includegraphics[width=0.9\linewidth]{fig2_latency_gap2.png}
\caption{Board latencies (inference vs inference + update) against the Gap~2 budget
(\SI{100}{\milli\second}). \emph{Measured on-board} (\texttt{exp\_S36\_summary.json}).}
\label{fig:latency}
\end{figure}

\subsection{Collapse of naive PTQ (measured on-board)}

The \emph{legacy} INT8 quantization of the binary EWC head collapses on the board: F1 drops from
$\approx$\,\num{0.92} (FP32) to \num{0.138} (Pronostia) and \num{0.1337} (Monitoring), for an
INT8$\leftrightarrow$FP32 agreement of only $\approx$\,\numrange{0.60}{0.74}. Weight RAM is indeed divided by
4, but the metric is destroyed (Table~\ref{tab:int8board}).

\begin{table}[t]
\centering
\caption{INT8 \emph{legacy} vs FP32 on-board (\emph{measured}, \texttt{exp\_S36}). RAM $\div 4$ but collapsed
F1.}
\label{tab:int8board}
\begin{tabular}{lcccc}
\toprule
Dataset & F1 FP32 & F1 INT8 legacy & INT8/FP32 agr. & RAM $\div$ \\
\midrule
Pronostia  & \num{0.9164} & \num{0.138}  & \num{0.7364} & $\times 4$ \\
Monitoring & \num{0.9194} & \num{0.1337} & \num{0.5955} & $\times 4$ \\
\bottomrule
\end{tabular}
\end{table}

\subsection{Diagnosis and recovery (bit-exact PC emulation)}

A \emph{bit-exact PC-emulated} ablation (exact kernel reference) isolates the cause and demonstrates recovery
(Table~\ref{tab:ablation}, Figure~\ref{fig:ablation}). The dominant factor is the \emph{calibrated
per-tensor scale} on an \texttt{int32} accumulator: \texttt{per\_tensor\_calib} restores $+\,\num{0.88}$ F1
(Pronostia \num{0.9462}, Monitoring \num{0.9201}). The \texttt{q15} format exactly matches FP32.

\begin{table}[t]
\centering
\caption{INT8 ablation ladder (F1, \emph{bit-exact PC emulation}, \texttt{exp\_S39\_ablation}).}
\label{tab:ablation}
\begin{tabular}{lcc}
\toprule
Scheme & Pronostia F1 & Monitoring F1 \\
\midrule
FP32 (ref.)        & \num{0.9616} & \num{0.9194} \\
\texttt{legacy\_c} & \num{0.0663} & \num{0.1178} \\
\texttt{per\_tensor\_calib} & \textbf{\num{0.9462}} & \textbf{\num{0.9201}} \\
\texttt{per\_channel\_int8} & \num{0.9426} & \num{0.9187} \\
\texttt{q15}       & \num{0.9616} & \num{0.9194} \\
\bottomrule
\end{tabular}
\end{table}

\begin{figure}[t]
\centering
\includegraphics[width=0.9\linewidth]{fig3_ablation_ladder.png}
\caption{Ablation ladder \texttt{legacy\_c}$\to$\texttt{per\_tensor\_calib}$\to$\texttt{per\_channel}
$\to$\texttt{q15}. \emph{Bit-exact PC emulation} (\texttt{exp\_S39\_ablation}).}
\label{fig:ablation}
\end{figure}

\begin{figure}[t]
\centering
\includegraphics[width=0.9\linewidth]{fig4_int8_recovery_board.png}
\caption{INT8 vs FP32: \emph{legacy} collapse then recovery. Collapse and recovery are both
\emph{measured on-board} (Sprints~36 and~40); the \emph{PC-emulated} per-channel column (Sprint~39) is
recalled with hatching. The \texttt{q15} format on-board remains "to be measured".}
\label{fig:recovery}
\end{figure}

\subsection{Recovery measured on-board (calibrated v2 kernel)}
\label{sec:recovery-board}

Recovery is no longer an emulator prediction: the calibrated-scale v2 kernel was flashed onto the
NUCLEO-F439ZI and streamed over the four per-channel cells (two datasets $\times$ two protocols). F1 returns
to the FP32 level --- \num{0.9173} against \num{0.9194} in FP32 on Monitoring, \num{0.8995} against
\num{0.9164} on Pronostia --- while weight RAM stays divided by 4. Parity against the bit-exact emulator is
\emph{exact} in the frozen protocol (\num{1.000}), which validates a lossless port, and no CRC error was
observed (Table~\ref{tab:recovery}).

\begin{table}[t]
\centering
\caption{INT8 recovery \emph{measured on-board} (per-channel v2 kernel, \texttt{exp\_S40\_board\_v2}).
Parity = agreement with the bit-exact emulator; agreement = INT8 against FP32 on the same samples.}
\label{tab:recovery}
\begin{tabular}{llccccc}
\toprule
Dataset & Protocol & F1 & Parity & Agreement & Latency (\micro s) & \texttt{.bss} (B) \\
\midrule
Monitoring & frozen & \num{0.9173} & \num{1.000}  & \num{0.9996} & 65  & 101236 \\
Monitoring & online & \num{0.9016} & \num{0.9885} & \num{0.8337} & 577 & 101236 \\
Pronostia  & frozen & \num{0.8995} & \num{1.000}  & \num{0.9951} & 68  & 106152 \\
Pronostia  & online & \num{0.9212} & \num{0.9736} & \num{0.8100} & 606 & 106152 \\
\bottomrule
\end{tabular}
\end{table}

Two caveats deserve to be stated. First, the INT8$\leftrightarrow$FP32 agreement drops online
(\num{0.8337} and \num{0.8100} against $>\,\num{0.99}$ frozen): the continual update makes the two weight
trajectories diverge, without the final metric suffering from it. Second, the \texttt{q15} format and the
A/B rerun of the \texttt{int8\_legacy} scheme under the v2 kernel were not streamed: those cells remain
"to be measured", in line with the "no invented number" rule.

\subsection{RAM $\times$ F1 $\times$ latency trade-off}

Figure~\ref{fig:pareto} synthesizes the trade-off: the calibrated kernel brings the metric back to FP32 level
while keeping the RAM gain (INT8 $\div 4$, Q15 $\div 2$), with no prohibitive latency cost. The star marks the
point actually flashed (per-channel v2 kernel); the other points are emulated. The expected latency gain,
however, does not materialize --- we return to this in Section~\ref{sec:discussion}.

\begin{figure}[t]
\centering
\includegraphics[width=0.9\linewidth]{fig5_pareto_ram_f1_latency.png}
\caption{RAM $\times$ F1 $\times$ latency trade-off front. Mixed \emph{measured}/\emph{emulated} sources
annotated per point.}
\label{fig:pareto}
\end{figure}
